\subsection{Process and Outcome Supervision}
\label{sec:4.2.5}
Everything above rewards a system for what it produced. Process supervision rewards it for how it got there, scoring each step of a chain of reasoning rather than the answer at the end. The technique became possible at scale once models were prompted to write the intermediate steps out — Jason Wei and coauthors showed in 2022 that asking for a series of intermediate steps sharply improves performance on multi-step problems \autocite{wei2022chain} — because a step that is written down is a step that can be marked.

The comparison has been run directly. Jonathan Uesato and coauthors, working on math word problems in 2022, found that outcome-based supervision matched process-based supervision on final-answer accuracy while using less labeled data, and diverged sharply on the reasoning underneath: among answers that came out right, the rate of faulty reasoning was 14.0 percent under outcome supervision and 3.4 percent under process supervision \autocite{uesato2022solving}. Hunter Lightman and coauthors reported in 2023 that on a harder benchmark process supervision beat outcome supervision on final answers as well, solving 78 percent of a representative subset of the MATH dataset \autocite{lightman2023lets}.

The first result is the one that bears on chapter~\ref{sec:3}. A system supervised on outcomes can be right for reasons that will not survive a change of circumstance, and nothing in the training signal distinguishes it from a system that was right for reasons that will. Chapter~\ref{sec:3} turns on exactly that distinction: a floor has to be a reason held rather than a result produced, because a result can be produced by machinery that will produce a different result under different pressure. Process supervision is the first stage in this pipeline that addresses the distinction at all, and it is evidence that the distinction is buildable and not only arguable.

Its limit is in the same place. What gets supervised is the reasoning the system states, and section~\ref{sec:11.1} carries the evidence that a model's report about its own processing corresponds to that processing partially, unreliably, and in a way that shifts with how the question is asked. A system rewarded for producing acceptable-looking reasoning has been given a reason to produce acceptable-looking reasoning. That is not an argument against the method, which is better than the alternative on the evidence above. It is the reason section~\ref{sec:3.7} declines to settle whether a bearer can be believed on the strength of what it says about itself, and it puts a ceiling on how much of that problem this stage can be asked to carry.
